\chapter{Proposed Solution}

\section{Problem Statement}

The RISC-V ISA is designed with intent to support custom extensions,
so as to allow developers to optimize processors for specific workloads.
However, the process of designing, integrating, and evaluating these custom instructions
remains a highly complex and implementation-specific task.

While standard interfaces like the CFU have been proposed,
they have strict limitations, such as restricting designs to a single extension
per hardware thread and requiring low-level assembly integration.

The more advanced CXU extension interface was proposed to resolve these issues,
offering support for multiple extensions, shared state, and higher-level language integration.

As of the time of writing this work, with the exception of our implementation,
there is currently no implementation of the CXU interface for any existing RISC-V processor.
As so, there is also no unified sandbox environment that allows developers
to prototype, test, and benchmark these custom hardware units
under realistic, operating system-level conditions.

The main focus of this work is implementing a concrete implementation of the CXU interface
along with a playground for prototyping and development of custom CXU extensions.
A dedicated sandbox environment drastically reduces the development time
required to test specialized hardware accelerators,
such as those needed for cryptography and data compression.

Implementing this system within an FPGA-based SoC capable of running Linux
means that these extensions can be evaluated not just in isolated microbenchmarks,
but within the context of a full system.

Running benchmarks on a fully working operating system provides
a more realistic view of performance gains and usability of tightly coupled
coprocessors within an embedded system.

\section{System Architecture and Playground Design}

The proposed sandbox is built upon the litex framework.
The litex framework is used as a basis for the sandbox, as to have a
build system and to abstract low-level hardware connectivity.
Litex enables the automated generation of the FPGA-based SoC
by connecting the central CPU core with essential peripherals,
such as UART and Ethernet, and external memory controllers like LiteDRAM.
Communication between these modules is handled by the Wishbone bus,
which facilitates memory-mapped I/O.

The core processor utilized in this SoC is VexiiRiscv,
an RV32IMA configured CPU implemented in the SpinalHDL hardware description language.
VexiiRiscv's plugin-based architecture is highly parameterized,
allowing for scalable feature composition and the clean integration
of custom execution units.

The entire playground is available as a single repository \cite{CXUPlayground}.
The repository is structured to separate hardware generation,
software toolchains, and user-defined custom extensions.
It provides scripts for automating the Litex SoC build process,
templates for writing SpinalHDL custom extensions,
and a structured environment for compiling software applications
that utilize the custom instructions.
\section{CXU Integration in VexiiRiscv}

To provide low-latency instruction-level acceleration,
the CXU mechanism is integrated as a tightly coupled accelerator
directly within the CPU pipeline.
Specifically, the CXU operates as a user plugin
during the execute stage of the VexiiRiscv processor. 

Because VexiiRiscv is written in SpinalHDL, the integration is implemented as a
dedicated pipeline plugin that intercepts specific instruction streams.
The data flow integration is designed as such:

\begin{itemize}
    \item Specific opcode patterns reserved for custom instructions
    are mapped to the CXU during the decode stage.
    \item The CXU command receives source operands directly
    from CPU registers \textbf{rs1} and \textbf{rs2}.
    \item The custom logic processes these inputs based on its internal state
    and returns the output to the destination register \textbf{rd}.
    \item Software manages and queries the CXU's shared state
    via dedicated \textbf{cxidx} (Extension Index) and \textbf{cxdata} Control/Status Registers (CSRs).
    \item For multi-cycle operations, the plugin asserts a halt signal to the CPU pipeline, managing pipeline
    stalls and structural hazards to ensure correct execution ordering and data integrity.
\end{itemize}

\begin{figure}[htb]
    \centering
    \begin{tikzpicture}[
        node distance=0.8cm and 0.5cm,
        stage/.style={draw, rectangle, minimum width=2cm, minimum height=1.2cm, fill=blue!10, rounded corners, font=\bfseries},
        cxu/.style={draw, align=center, rectangle, minimum width=3cm, minimum height=1.5cm, fill=green!10, rounded corners, font=\bfseries},
        arrow/.style={-Latex, thick},
        signal/.style={-Latex, dashed, thick, red}
    ]
    \node[stage] (fetch) {Fetch};
    \node[stage, right=of fetch] (decode) {Decode};
    \node[stage, right=of decode] (execute) {Execute};
    \node[stage, right=of execute] (memory) {Memory};
    \node[stage, right=of memory] (writeback) {Writeback};
    
    \node[cxu, below=1.5cm of execute] (cxu_unit) {CXU Plugin \\ \small (Custom Accelerator)};
    
    \draw[arrow] (fetch) -- (decode);
    \draw[arrow] (decode) -- (execute);
    \draw[arrow] (execute) -- (memory);
    \draw[arrow] (memory) -- (writeback);
    
    \draw[arrow] (decode.south) -- ++(0,-0.25) -| node[pos=0.75, left, font=\small] {Opcode Map} ([xshift=-1cm]cxu_unit.north);
    \draw[arrow] ([xshift=-0.4cm]execute.south) -- node[pos=0.5, above, font=\small] {rs1, rs2} ([xshift=-0.4cm]cxu_unit.north);
    \draw[arrow] ([xshift=0.4cm]cxu_unit.north) -- node[right, font=\small] {rd} ([xshift=0.4cm]execute.south);
    
    \draw[signal] (cxu_unit.east) -| node[pos=0.75, right, font=\small, text=black] {Halt / Pipeline Stall} (memory.south);
    
    \end{tikzpicture}
    \caption{Integration of the CXU as a tightly coupled plugin within the VexiiRiscv pipeline}
\end{figure}

\section{Sandbox Development Environment and Buildroot Integration}

The core contribution of this work is the unified sandbox
that allows users to seamlessly design, test, and benchmark CXU hardware.
Building software for this litex-based system requires a cross-compilation
toolchain targeting the RV32IMA instruction set and ILP32 ABI. 

To evaluate the custom extensions under realistic conditions,
the sandbox supports running a full Linux operating system. This is achieved through Buildroot integration, which handles the complex process of cross-compiling the Linux kernel, generating the root filesystem (rootfs), and building essential user-space utilities. 

Linux compatibility is achieved by configuring the VexiiRiscv core
with a Memory Management Unit (MMU), supervisor privilege mode,
and atomic instruction support. 
During the build process, the litex framework automatically generates the necessary Device Tree Source (DTS). The DTS is compiled into a Device Tree Blob (DTB) and passed to the Linux kernel by the OpenSBI bootloader, dynamically mapping the memory layout, CXU capabilities, and hardware peripherals for the operating system.

\begin{figure}[htb]
    \centering
    \begin{tikzpicture}[
        layer/.style={draw, rectangle, minimum width=8.5cm, minimum height=1.2cm, font=\bfseries},
        hw/.style={fill=gray!20},
        fw/.style={fill=purple!20},
        os/.style={fill=orange!20},
        rt/.style={fill=cyan!10},
        usr/.style={fill=cyan!30},
        box/.style={rectangle, draw, rounded corners=3pt,
            minimum width=3.2cm, minimum height=0.75cm, align=center,
            text centered, font=\small}
    ]
    \node[box, hw] (hardware) {Hardware: LiteX SoC + VexiiRiscv + CXU Extensions};
    \node[box, fw, above=0cm of hardware] (firmware) {Firmware: OpenSBI \& Device Tree Blob (DTB)};
    \node[box, os, above=0cm of firmware] (linux) {Operating System: Linux Kernel (MMU enabled)};
    \node[box, rt, above=0cm of linux] (runtime) {CXU Runtime ABI \& C-level Library};
    \node[box, usr, above=0cm of runtime] (apps) {User Space Applications (AES, GZIP Benchmarks)};
    
    \node[left=0.8cm of hardware, align=right, font=\itshape] {Kernel Mode};
    \node[left=0.8cm of runtime, align=right, font=\itshape] {User Mode};
    
    \end{tikzpicture}
    \caption{The CXU Sandbox full-system stack}
\end{figure}

\section{The CXU Runtime}

A critical component of bringing CXU extensions to an operating system environment is the CXU runtime software interface. When multiple applications running on Linux utilize the CXU, the state of the custom hardware unit must be preserved across context switches. 

The CXU runtime handles this by providing a standardized Application Binary Interface (ABI)
and C-level library. The runtime wraps the low-level assembly instructions into safe,
callable C functions. It manages the \textbf{cxidx}
register to securely multiplex access to different hardware extensions
among various software threads. 

Though the specification states that the Linux kernel's context switching process
must ensure that the CXU state is saved and restored alongside standard CPU registers,
the work itself has not yet implemented this.

\section{Validation through Case Studies}

To validate the effectiveness and flexibility of the CXU sandbox,
the framework will host and benchmark complex cryptographic
and data compression algorithms. 

\begin{itemize}
    \item \textbf{AES Encryption Optimization:} Advanced Encryption Standard workloads
    will be accelerated by offloading computationally expensive finite field
    operations in $GF(2^8)$ to the CXU.
    Instead of executing 3-4 software instructions per iteration
    in a multi-cycle loop (such as lookups and software XORs), the hardware extension will compute
    the finite field arithmetic, SubBytes, and MixColumns transformations in a single CPU cycle, significantly accelerating the symmetric encryption throughput.
    
    \item \textbf{GZIP Compression Optimization:} The DEFLATE algorithm relies heavily
    on pattern detection and entropy encoding.
    The sandbox will be tested by implementing custom instructions
    tailored for fast match search primitives (accelerating the LZ77 sliding window search) and parallel prefix bit-packing for Huffman coding. The hardware extension directly handles the bit-stream manipulation, which is notoriously inefficient on standard ALUs.
\end{itemize}

These case studies will serve as a proof of concept,
demonstrating how modular hardware extensions can yield substantial
runtime improvements and cycle count reductions
within the developed framework.

\section{Testing and Benchmarking}

An extensive testing methodology is established within the playground
to measure the correctness and performance of the implemented extensions.
The testing strategy is twofold:

\begin{itemize}
    \item \textbf{Unit Testing and Verification:} Each CXU extension is initially simulated in a
    Verilator-based testbench to ensure logical correctness before synthesis.
    Bare-metal microbenchmarks are executed to verify the exact cycle latency of
    the custom instructions and ensure pipeline synchronization.
    
    \item \textbf{System-Level Benchmarking:} After beins synthesizes on the FPGA,
    benchmarks are also executed within the Buildroot Linux environment.
    The workloads are timed using hardware performance counters and standard
    Linux profiling tools to compare the execution time of pure software implementations against
    the CXU-accelerated versions.
\end{itemize}

\section{General Methodology}

The data collection process for this work is constructed of
phases: the implementation of the CXU interface itself,
and the implementation of concrete hardware extensions used as experimental subjects.

The first part is implementing the CXU extension mechanism within
the VexiiRiscv pipeline. This required implenenting a
dedicated SpinalHDL plugin that intercepts
custom-opcode instructions during the execute stage,
routes operands from registers \textbf{rs1} and \textbf{rs2} to the coprocessor,
manages pipeline stalls for multi-cycle operations, and writes results back to \textbf{rd}.
The \textbf{cxidx} and \textbf{cxdata} Control and Status
Registers (CSRs) were introduced to allow software to select an active
extension and exchange shared state.

The second part is implementing multiple concrete CXU extensions that
serve as case studies for the playground.

The first extension is an AES extension offloads finite
field arithmetic in $\text{GF}(2^8)$, specifically the \textbf{SubBytes} and
\textbf{MixColumns} transformations, to dedicated hardware logic, replacing
multi-instruction software loops with a single custom instruction per
operation. Second is the GZIP extension. There are multiple versions of these
extensions, one for a basic arithmetic optimization, and one for the full
GZIP decompresison extension.

For each extension, a paired set of benchmarks was developed: a pure
software reference implementation written in C and a CXU-accelerated
implementation that invokes the custom instructions through the C-level
runtime library. Both variants were compiled for the \textbf{RV32IMA}
target using the same cross-compilation toolchain and optimisation flags to
ensure a fair comparison. Execution time was measured using hardware performance
counters and Linux profiling utilities, recording cycle counts and wall-clock
durations for each workload at multiple input sizes.

\section{Data Analysis}

Performance analysis consists of comparing the cycle counts and execution
times recorded for the software-only reference implementations against those
of the CXU-accelerated variants for the same inputs.

Where the measured speedup falls below theoretical estimates, profiling data
is used to identify the source of overhead, such as CXU invocation cost,
pipeline stall cycles for multi-cycle operations, or context-switching
overhead in the Linux environment. Results from bare-metal microbenchmarks
are contrasted with full-system Linux results to assess the overhead
introduced by the operating system and the CXU runtime library.

\section{Limitations}

Several limitations of the current implementation affect the practical
utility of certain CXU features and must be taken into account when
interpreting the benchmarking results.

The first and most significant limitation concerns the shared-state mechanism
provided by the \textbf{cxidx} and \textbf{cxdata} CSRs.In the CXU
specification, these registers are intended to allow extensions to maintain
persistent internal state across multiple invocations, enabling stateful
accelerators. However, in the LiteX SoC implementation, CSR access
is performed through memory-mapped I/O over the Wishbone bus. Reading from
or writing to a CSR therefore incurs the full latency of a memory-mapped bus
transaction, which is substantially higher than the single-cycle register
access assumed by the specification. As a consequence, extensions that rely
heavily on reading or writing shared state through \textbf{cxdata} between
invocations will experience significant overhead, effectively negating any
computational savings provided by the custom hardware logic. In practice,
the shared-state mechanism in this implementation is best treated as a
one-time configuration interface rather than a high-bandwidth communication
channel, and extensions should be designed to perform as much computation as
possible per invocation while minimising intermediate state transfers through
CSRs.

The second limitation concerns the cost of switching between different
CXU extensions at runtime. The CXU specification supports up to
$2^{\text{XLEN}}$ distinct extensions, with the active extension selected by
writing the desired index to the \textbf{cxidx} CSR. Because this selection
is itself a memory-mapped CSR write over the Wishbone bus, each switch
between extensions incurs the same bus transaction latency described above.
Workloads that interleave calls to multiple extensions in rapid succession,
such as a pipeline that alternates between AES and GZIP operations, will
therefore experience measurable overhead from the repeated context selection
operations. To achieve the best possible performance, software should be
structured to batch all invocations of a given extension together before
switching to another, minimising the number of \textbf{cxidx} writes
relative to the number of productive custom instruction calls.

Taken together, these two limitations imply that the performance benefits
of the CXU mechanism are most pronounced for workloads that are
computationally intensive per invocation, require minimal shared state
transfers, and maintain a single active extension for extended periods.
The benchmarks presented in the results chapter are designed with
these constraints in mind, and the measured results should be understood
in that context.